\section{Experimental Setup}

\subsection{Evaluation Questions and Protocol}

The experiments address seven questions: whether peptide sequence retains tissue-associated signal after matching the \plainMHC{} and source protein; the performance of the human model; the contribution of major model components; the effect of removing every held-out peptide from every training task; replication of task learnability in mouse; the signal recoverable without tissue identity; and whether differences between model structures remain under \plainPeptideOOF{}. These questions connect the technical comparisons to their biological interpretation. The human--mouse analysis concerns replication of the task definition, not transfer of one model between species.

Many model choices were explored on the ordinary human benchmark. To avoid continuing to adjust the analysis after seeing the harder results, we specified the peptide-separated evaluation in advance. This plan was \plainFrozen{} and specified the data split, model versions, \plainSeeds{}, training budget, comparison models, performance measures, and statistical tests. The exact split files, configurations, and timestamps are saved with the predictions.

The two main biological questions should not be treated as interchangeable. Evaluation on familiar tasks with cross-task peptide reuse allowed (standard evaluation) asks whether the model can rank held-out pairs within familiar tissue--antigen-presenting-molecule combinations (tissue--MHC combinations), even though an individual peptide may have appeared in another training task. \plainPeptideOOF{} asks the harder question of whether ranking remains possible for entirely new peptide identities in those same combinations. The first describes the internal benchmark; the second measures robustness to peptide reuse. Table~\ref{tab:evaluation-protocols} defines the data pools, overlap rules, held-out-fold cross-validation status (OOF status), matching, and model repeats.

For every trainable method, held-out examples are excluded from parameter fitting within that run. In \plainPeptideOOF{}, a held-out peptide and every pair connected to that peptide remain outside the corresponding fitting fold. Thus, a reported prediction cannot benefit from direct exposure to the same peptide identity during that fold's training. This within-run separation does not imply that the overall model design was chosen without earlier access to the ordinary benchmark; the selected structures had already been informed by \plainPairOOF{} on the same training pool.

To assess sensitivity to removing peptide reuse, we also create an ordinary split from exactly the same pool of pairs and with nearly the same number of held-out pairs per task (\plainMatchedOOF{}). The intended design difference between this same-pool reference and \plainPeptideOOF{} is whether an exact peptide may cross the training boundary. However, grouping pairs linked directly or indirectly by a shared peptide can also change fold composition, including parent-protein and related-motif distributions. We therefore interpret the observed difference as a robustness gap under peptide separation, not as a pure causal estimate of the effect of repeated peptide identity.

\subsection{Training and Baseline Models}

All sequence-to-vector components and task-specific output layers (encoders and heads) are optimized with decoupled weight-decay optimization (AdamW) \cite{loshchilov2019adamw} for 25 complete passes through the training schedule (epochs), using 512 rows per parameter update (batch size 512), a learning rate of \(10^{-3}\), and weight decay of \(10^{-4}\). The size of each parameter update is capped at 1.0 (gradient-norm clipping). The pathway shared across tasks and the pathway restricted to one \plainHLA{} allele (global and HLA-specific branches) are trained independently for each \plainSeed{}. Within every \plainOOF{} experiment, only the fitting partition is used to estimate model parameters; held-out rows are used solely for prediction and evaluation.

The human main model uses three \plainSeeds{}: 20260704, 20260705, and 20260706. The model structure (architecture), parameter-update settings (optimizer settings), weights for the additional training tasks (auxiliary weights), score-combination rule (fusion rule), and set of \plainSeeds{} are \plainFrozen{} for the reported evaluation. Because the training schedule uses a fixed number of complete passes (epochs), the held-out fold is not used to decide when to stop training (early stopping). The same preprocessing and task mappings are reconstructed from each fitting partition so that no held-out labels influence training.

The comparison models range from simple sequence rules to models that share information across tissue--antigen-presenting-molecule combinations (tasks). A linear model using a separate indicator for each amino acid and position (one-hot logistic regression) asks what can be learned from amino-acid identity with a linear rule. A collection of decision trees using established amino-acid similarities (\plainBLOSUM{} random forest) adds known similarity between substitutions. The \plainSharedHeads{} learns one common peptide representation with a separate output for each tissue--antigen-presenting-molecule combination (tissue--MHC combination). Models with one broadly shared pathway and one allele-restricted pathway (dual-branch models) additionally compare broad sharing with sharing restricted to one \plainHLA{} allele. All methods within a table use the same set of tissue--antigen-presenting-molecule combinations (task inventory) and data split, so differences are not created by evaluating different tissues or alleles.

The collection of decision trees using \plainBLOSUM{} features \cite{henikoff1992blosum} joins the 20 substitution scores at each of the nine peptide positions into 180 input values and fits a separate model per task. It uses 200 trees, tests a square-root-sized random subset of input values at each split (square-root feature subsampling), and requires at least two examples in a terminal tree group (minimum leaf size two). The two-pathway feed-forward controls (MLP dual-branch baselines) use 16-number residue vectors followed by joining all positions (flattening) and two 128-unit layers in a \plainMLP{}, with a rectified linear unit (ReLU) and 20\% random omission during training (dropout 0.2) after the first layer. The \plainAuxDual{} adds tissue and \plainHLA{} training-guidance losses weighted 0.1 each to its across-task pathway; the \plainPlainDual{} removes those additional objectives while retaining the two pathways and their fixed equal-weight combination of within-task orderings (task-rank fusion). Human neural-model averages use three \plainSeeds{} (20260704--20260706). Mouse results under \plainPeptideOOF{} average five \plainSeeds{} (20260704--20260708).

Every model structure evaluated under \plainPeptideOOF{} reads the saved file that assigns linked groups to folds (component manifest) rather than regenerating folds. The human comparison contains 73,798 pairs in 17,036 groups linked directly or indirectly by shared peptides (connected components) across 157 tasks; the mouse comparison contains 6,766 pairs in 1,844 such groups across 24 tasks. Both use three evaluation folds (outer folds) with zero complete-pair and exact-peptide overlap between fitting and held-out partitions. Using identical folds for every model separates model-structure comparisons from variation caused by regenerating the split.

\subsection{Controls Without Tissue Information}

The most direct biological control removes the tissue label while retaining peptide and \plainMHC{} information (\plainMHCOnly{}). If this control performed as well as the full model, the benchmark could be explained largely by general compatibility between a peptide and its antigen-presenting molecule (peptide--MHC compatibility). The human and mouse \plainMHCOnly{} models use nearly the same model capacity and exactly the same evaluation splits as their full counterparts; the key information removed is tissue identity.

We also score every peptide--antigen-presenting-molecule query (peptide--MHC query) with MHCflurry and NetMHCpan, two established predictors that do not use the target tissue \cite{odonnell2020mhcflurry,reynisson2020netmhcpan}. These \plainGeneralPMHC{} quantify how much of the tissue benchmark constructed from \plainMatchedPairs{} can be explained by general binding or presentation propensity. They are not retrained on the TissuePMHC labels.

The control set contains 79,759 unique human peptide--antigen-presenting-molecule queries (peptide--MHC queries) across 35 \plainHLA{} alleles and 6,663 unique mouse queries across four \plainHtwo{} restrictions. The predicted presentation score from MHCflurry, the \plainEL{} from NetMHCpan, and the \plainBA{} from NetMHCpan all cover 100\% of rows, pairs, and tasks, so comparisons use exactly the same evaluable pairs as the tissue-conditioned models. Controls are evaluated on the \plainFixedTest{} and on the complete training pools used for \plainMatchedOOF{} and \plainPeptideOOF{}. Because the external scores were \plainFrozen{} rather than fitted within each fold, their standalone score metrics are identical for the two \plainOOF{} partitions of the same pool; protocol-specific comparisons use the corresponding tissue-conditioned \plainOOF{} predictions.

Finally, we test whether adding a tissue-blind general presentation score supplies information missing from the tissue-conditioned score. The rule combining the two scores is learned from other folds and then applied to the held-out fold (cross-fitted stacking), so it is learned without using the rows on which it is evaluated. This is a diagnostic analysis, not a newly proposed final model.

\subsection{Performance Metrics and Statistical Analysis}

The main question is whether reported target-tissue peptides tend to rank above comparison peptides sharing presenting restriction and parent protein (matched comparison peptides). AUROC summarizes this ordering within each tissue--antigen-presenting-molecule task (tissue--MHC task). We calculate it separately for every task and then average tasks equally, so tissues or alleles with more rows do not dominate:
\[
\operatorname{MacroAUROC}
=
\frac{1}{|\mathcal{T}_{\mathrm{valid}}|}
\sum_{\tau\in\mathcal{T}_{\mathrm{valid}}}
\operatorname{AUROC}_{\tau},
\]
where \(\mathcal{T}_{\mathrm{valid}}\) contains the prespecified tasks with both labels available for evaluation. Equal-weight averaging across tasks (task-macro averaging) gives every task equal weight regardless of its sample size. The equal-weight average of AUPRC across tasks (task-macro AUPRC) is defined analogously:
\[
\operatorname{MacroAUPRC}
=
\frac{1}{|\mathcal{T}_{\mathrm{valid}}|}
\sum_{\tau\in\mathcal{T}_{\mathrm{valid}}}
\operatorname{AUPRC}_{\tau}.
\]
Because the paired dataset has a constructed 1:1 label ratio, AUPRC is interpreted relative to the per-task benchmark prevalence and is not compared directly with AUPRC from naturally imbalanced clinical cohorts.

An average can hide weak performance in difficult tissues, so we also report the lowest-performing ten human tasks and six mouse tasks. \plainPairAcc{} gives the most direct interpretation: it is the fraction of original \plainMatchedPairs{} for which the recorded-positive peptide receives the higher score. Accuracy, F1, and MCC require a fixed cutoff and are secondary because the final rank score is not a calibrated biological probability.

Metrics are computed in the following order. For each \plainSeed{}, predictions from all held-out evaluation folds (outer held-out folds) are first joined. Each task metric is then computed from its complete \plainOOF{} prediction vector, after which task metrics are averaged with equal weight. Averaged-model predictions under \plainPeptideOOF{} are produced only from model members trained on the corresponding fitting fold. Metrics calculated separately within each fold (fold-level metrics) are reported descriptively but are not averaged as if folds were independent experimental replicates. Means and standard deviations across \plainSeeds{} quantify training variation, whereas the average across models (ensemble) remains a single final predictor.

Uncertainty intervals are obtained by repeatedly resampling complete tissue--antigen-presenting-molecule tasks (tissue--MHC tasks) rather than treating folds or repeated model runs as independent biological samples. Paired tests compare models task by task, and multiple comparisons are corrected within each planned analysis family. These statistics quantify consistency across the recorded tasks, but the tasks are not fully independent: several may share a tissue, an allele, proteins, or related peptides. We therefore emphasize effect size, the number of tasks improving or declining, and agreement between species rather than treating the results as evidence from independent patient cohorts.

\subsection{Cross-species Interpretation and Reproducibility}

Human and mouse use the same definition of a peptide pair sharing presenting restriction and parent protein (matched-pair definition), minimum number of pairs required for each biological combination (minimum task size), \plainFixedTest{} design, \plainPairOOF{} design, and \plainPeptideOOF{} design. They differ in the number of tissue--antigen-presenting-molecule combinations (tasks), \plainMHC{} system, data composition, and selected model. Concordant findings across species therefore support replication of whether the benchmark question can be learned (task learnability), but do not establish a species-independent molecular mechanism or a controlled comparison of model structures (architecture comparison).

The human model uses three \plainSeeds{}: 20260704, 20260705, and 20260706. The mouse averaged model (ensemble) uses those three \plainSeeds{} plus 20260707 and 20260708. Fold assignments for \plainPairOOF{} and for groups linked by shared peptides in \plainPeptideOOF{} use the recorded \plainSplitSeed{} 20260711. Comparisons across species emphasize whether both benchmarks retain above-random signal and whether structured information sharing improves upon simple sharing under the same protocols. We do not compare the numerical advantage of the two model structures (architectures) as a controlled species effect, because data size, task diversity, and model selection differ.

For each experiment, we retain the exact rules governing access to fitting and test data (train/test access policy), the file assigning rows or connected groups to folds (split manifest), the \plainSplitSeed{} used to create folds, the \plainSeeds{} used to train models, model configuration, and command-line arguments. Reports include the number of complete training passes (epoch count), rows per update (batch size), parameter-update method (optimizer), learning rate, weight decay, random omission rate (dropout), parameter count, training time, peak memory when available, hardware, Python and PyTorch versions, and handling of \plainSeeds{}. Predictions are saved at the individual-row level with tissue--antigen-presenting-molecule combination (task), label, pair identifier, fold, model member, and final score whenever supported by the run.

Dataset and configuration files are accompanied by short digital fingerprints used to detect file changes (hashes), and each result directory records the timestamped protocol and code state. The metadata for the mouse \plainFixedTest{} explicitly record that test data were accessed only for evaluation of the model that had been \plainFrozen{} and that model selection was not performed on the test set. Runs under \plainPeptideOOF{} retain their assignments of peptide-linked groups (component assignments) and checks confirming zero forbidden overlap (zero-overlap audits).

\subsection{Data and Code Availability}

The accompanying project repository contains the processed benchmark tables, files assigning pairs and peptide-linked groups to folds (pair and component split manifests), row-level predictions used for the reported analyses, configuration files, digital file fingerprints (hashes), and analysis scripts. Raw material derived from the \plainIEDB{} and third-party predictor assets remain subject to their original access and licensing terms; NetMHCpan executables are not redistributed.
